\section{Introduction}
\label{sec:introduction}
%% ============================================================

Any program of large-scale space infrastructure construction---whether solar power satellites, orbital habitats, or distributed energy collection architectures---eventually reaches a decision point that no amount of launch cost reduction can fully avoid under Earth-supplied logistics: at what scale does it become cheaper to manufacture components in space rather than to build them on Earth and haul them up through the gravity well? At megascale, some fraction of each unit's mass must still be Earth-sourced (hereafter ``vitamin'' components, \S\ref{sec:vitamin}). This question has been posed qualitatively since at least the 1970s \cite{oneill1974,oneill1976}, and it has recently gained renewed urgency as launch costs decline precipitously \cite{jones2018}, lunar ISRU programs advance under the Artemis architecture \cite{sanders2015}, and private ventures pursue asteroid resource extraction \cite{elvis2012,andrews2015}.

The intuition is rooted in a cost asymmetry that launch cost reduction cannot fully eliminate. Reusable vehicles have brought per-kilogram costs down by roughly two orders of magnitude, from \$20{,}000/kg to projections below \$500/kg for Starship-class vehicles. But per-kilogram launch costs decline more slowly than manufacturing costs: even under aggressive launch learning (LR$_L = 0.90$), the propellant and operations component ($\sim$\$200/kg) constitutes an assumed operational asymptote under Earth-supplied logistics that the learnable component cannot breach. ISRU inverts this cost structure: large fixed capital costs up front, but each additional unit produced from local materials costs only energy and labor. The fixed costs divide across a growing denominator; marginal costs decline with experience.

Despite the intuitive appeal of this argument, a quantitative gap persists in the literature. Existing ISRU economic analyses are overwhelmingly mission-specific: they evaluate the cost-effectiveness of producing lunar oxygen for propellant \cite{sanders2015}, extracting water ice for life support \cite{zubrin1996}, or mining asteroidal platinum-group metals for terrestrial markets \cite{sonter1997,elvis2012}. We are not aware of prior work that combines schedule-aware NPV crossover analysis with systematic uncertainty characterization for generic manufactured products.

A second gap concerns the treatment of manufacturing learning. The Wright learning curve \cite{wright1936}, well established in aerospace production economics \cite{wertz2011}, predicts that unit manufacturing cost decreases as a power function of cumulative production volume. This effect has been extensively documented for Earth-based aerospace systems \cite{argote1990,nagy2013} but we are not aware of it being combined with NPV timing analysis in a comparative Earth-versus-ISRU cost model for generic structural units. The omission matters because the two pathways exhibit fundamentally different learning dynamics: Earth manufacturing benefits from learning in fabrication but exhibits limited learning in launch (baseline $\mathrm{LR}_L = 0.97$; we sweep from 90\% to 100\% in \S\ref{sec:sensitivity}), while ISRU benefits from learning in both extraction and fabrication, albeit from a higher initial cost base and with a slower initial ramp-up. The long-run consequences of this asymmetry are not obvious without a quantitative model.

Our analysis is primarily relevant to non-revenue infrastructure (habitats, depots, structural frameworks) where cost minimization is the decision criterion. For revenue-generating systems such as space solar power, the opportunity cost of ISRU's deployment delay may dominate the cost savings (\S\ref{sec:discussion}), making the Earth pathway preferred even above the cost crossover.

This paper makes three contributions. First, we present a parametric cost model with net present value discounting and pathway-specific delivery schedules that expresses the total cost of producing $N$ units via the Earth-launch pathway and the ISRU pathway as explicit functions of production volume, with learning curves, an integrated production ramp-up schedule, ISRU mass penalty and transport costs, and financing costs applied to the appropriate components (\S\ref{sec:model}). Second, we embed this model in a 10{,}000-run Monte Carlo framework with correlated sampling that propagates uncertainty in eighteen independent stochastic parameters (plus two derived)---including the ISRU operational cost floor, production rate, and transport duration---at each of several fixed discount rates, yielding probability distributions, convergence analysis, and global sensitivity rankings rather than point estimates (\S\ref{sec:results}). Third, we propose a phased hybrid transition strategy and discuss its implications for near-term space policy (\S\ref{sec:discussion}).

%% ============================================================
